\section{时间反演与对称性破缺}
\label{chap:time-reversal-breaking}

空间对称由酉算符实现，时间反演却必须同时复共轭振幅。这个反幺正性使 $\Theta^2=-1$ 的系统出现 Kramers 成对，而普通特征标理论不能不加修改地直接套用。

对称性还可能由哈密顿量保留、却被所选基态破坏。此时群仍然限制简并极小值和序参量怎样变换，但物理态只保留一个稳定子群。把时间反演和自发破缺分开处理，可以避免把反幺正结构与普通群--子群降阶混为一谈。

\subsection{时间反演对称性}
\label{sec:time-reversal}

前面所有对称性都由酉算符实现。时间反演不同：它不是普通的酉空间变换，而是\textbf{反酉}变换。这个差别看似技术性，实际上决定了 Kramers 简并以及拓扑能带理论的许多基本结构。因此必须从 $\vb k\to-\vb k$、自旋翻转与复共轭开始，并明确区分 $\Theta^2=+1$ 与 $\Theta^2=-1$。


\begin{definitionplain}[时间反演算符]
时间反演算符记为 $\Theta$。它是反线性的：
\[
\Theta(a\ket\psi+b\ket\phi)
=a^*\Theta\ket\psi+b^*\Theta\ket\phi,
\]
并保持概率内积的模，满足
\[
\braket{\Theta\psi|\Theta\phi}
=\braket{\phi|\psi}.
\]
若
\[
\Theta\hat H\Theta^{-1}=\hat H,
\]
则称系统具有时间反演对称性。
\label{def:time-reversal}
\end{definitionplain}

时间反演下，经典对应量的变换为
\[
 t\to-t,
\qquad
\vb r\to\vb r,
\qquad
\vb p\to-\vb p,
\qquad
\vb L\to-\vb L,
\qquad
\vb S\to-\vb S.
\]
因此无磁场、无磁序的普通动能与标量势通常保持不变，而 Zeeman 项
\[
-\boldsymbol\mu\cdot\vb B
\]
在只翻转粒子自由度而固定外磁场时会破坏时间反演对称性。

对于无自旋粒子，在位置表象且哈密顿量可以取实的常见情形中，可以选
\[
\Theta=K,
\]
其中 $K$ 是复共轭。对自旋 $1/2$，一个标准相位约定是
\[
\Theta=-\ii\sigma_yK,
\label{eq:spin-half-time-reversal}
\]
并且
\[
\Theta^2=-1.
\label{eq:theta-square-minus}
\]
这个负号正是 Kramers 定理的来源。

时间反演的反酉性并不是人为规定，而是由量子力学正则对易关系与动量反向共同逼出来的。

可以从最基本的正则对易关系看出这一点。时间反演要求
\[
\Theta\hat{\vb r}\Theta^{-1}=\hat{\vb r},
\qquad
\Theta\hat{\vb p}\Theta^{-1}=-\hat{\vb p}.
\]
而量子力学有
\[
[\hat r_i,\hat p_j]=\ii\hbar\delta_{ij}.
\]
若 $\Theta$ 是普通线性酉算符，则共轭后左边变成
\[
[\hat r_i,-\hat p_j]=-\ii\hbar\delta_{ij},
\]
但右边的数 $\ii\hbar$ 不会改变，产生矛盾。若 $\Theta$ 是反酉的，它会把复数 $\ii$ 共轭成 $-\ii$，于是两边同时变号，正则对易关系保持成立。因此反酉性不是人为约定，而是由“位置不变、动量反号”与量子对易关系共同强迫出来的。


在 $\sigma_z$ 基下，复共轭 $K$ 对 Pauli 矩阵的作用为
\[
K\sigma_xK^{-1}=\sigma_x,
\qquad
K\sigma_yK^{-1}=-\sigma_y,
\qquad
K\sigma_zK^{-1}=\sigma_z.
\]
单独的 $K$ 不能把三个自旋分量全部反号，因此还需乘一个酉矩阵。取绕 $y$ 轴转 $\pi$ 的自旋旋转
\[
U_y(\pi)=\exp\!\left(-\frac{\ii\pi}{2}\sigma_y\right)=-\ii\sigma_y,
\]
定义
\[
\Theta=U_y(\pi)K=-\ii\sigma_yK.
\]
直接计算可得
\[
\Theta\sigma_i\Theta^{-1}=-\sigma_i,
\qquad i=x,y,z.
\]
于是时间反演确实把自旋翻转。再次作用一次，
\[
\Theta^2=-1,
\]
说明半整数自旋与无自旋/整数自旋的时间反演结构有本质区别。

在周期晶体中，某些 $\vb k$ 与 $-\vb k$ 只差一个倒格矢；这些点在时间反演下映回自身，称为时间反演不变动量。

时间反演不变动量条件
\[
-\vb k=\vb k+\vb G
\]
等价于
\[
2\vb k=-\vb G.
\]
因此在 $d$ 维 Bravais 晶格中，每个倒格基方向上的系数只能取 $0$ 或 $1/2$，共有
\[
2^d
\]
个不等价时间反演不变动量。三维晶体因此有八个时间反演不变动量，二维有四个。一些时间反演不变动量在具体晶格中会被空间点群进一步分成不同的星，例如立方晶格里常见 $\Gamma$、若干 $X$、$L$ 型点。

Kramers 定理保证自旋 $1/2$ 时间反演体系在每个时间反演不变动量都至少二重简并。拓扑绝缘体理论进一步利用这些时间反演不变动量上的 Kramers 对如何在整个 Brillouin 区连接来定义拓扑不变量；那已经超出本课程，但其最基本的代数根源就是本节证明的 $\Theta^2=-1$。

\begin{theorem}[Kramers 定理]
若系统具有时间反演对称性，且在所讨论的态空间上
\[
\Theta^2=-1,
\]
则每个能量本征态至少二重简并。具体地，若
\[
\hat H\ket\psi=E\ket\psi,
\]
则 $\Theta\ket\psi$ 具有同样能量，并与 $\ket\psi$ 正交。
\label{thm:kramers}

\end{theorem}

\begin{proof}
由 $\Theta\hat H\Theta^{-1}=\hat H$，
\[
\hat H\Theta\ket\psi
=\Theta\hat H\ket\psi
=E\Theta\ket\psi,
\]
所以时间反演态能量相同。再用反酉内积恒等式，取 $\ket\phi=\ket\psi$、$\ket\chi=\Theta\ket\psi$：
\[
\braket{\Theta\psi|\Theta^2\psi}
=\braket{\Theta\psi|\psi}.
\]
由 $\Theta^2=-1$，左边又等于
\[
-\braket{\Theta\psi|\psi}.
\]
故
\[
\braket{\Theta\psi|\psi}=0.
\]
因此 $\ket\psi$ 与 $\Theta\ket\psi$ 是两个线性无关、同能量的本征态。
\end{proof}

证明中的关键不是“时间反演让自旋翻转”这句物理直觉，而是反酉性与 $\Theta^2=-1$ 的组合。设 $|\psi\rangle$ 是能量本征态。若 $[\Theta,\hat H]=0$，则 $\Theta|\psi\rangle$ 具有同样能量。接下来要证明这两个态不是同一个方向。

利用反酉算符满足
\[
\langle\Theta\phi|\Theta\psi\rangle=\langle\psi|\phi\rangle,
\]
取 $|\phi\rangle=|\psi\rangle$、并把第二个态再作用一次 $\Theta$，可得
\[
\langle\Theta\psi|\Theta^2\psi\rangle
=\langle\Theta\psi|\psi\rangle.
\]
若 $\Theta^2=-1$，左边又等于 $-\langle\Theta\psi|\psi\rangle$，所以该内积只能为零。于是 $|\psi\rangle$ 与 $\Theta|\psi\rangle$ 正交，必然是两个线性无关的同能态。


时间反演算符的构造完成后，现在分析它在晶体中的物理后果：$\vb k$ 和 $-\vb k$ 态的配对，以及复共轭表示的等价性判据。

时间反演把动量反号，因此 Bloch 态满足
\[
\Theta:\quad \vb k\longmapsto-\vb k.
\]
若系统有时间反演对称性，则能带满足
\[
E_n(\vb k)=E_n(-\vb k),
\label{eq:tr-band-symmetry}
\]
这里若存在自旋--轨道耦合，$n$ 应理解为连续追踪的 Kramers 相关能带，而不能简单把态写成固定的“$\uparrow/\downarrow$”自旋本征态。

若某个波矢满足
\[
-\vb k=\vb k+\vb G,
\]
则它称为时间反演不变动量（时间反演不变动量）。在自旋 $1/2$ 且 $\Theta^2=-1$ 的系统中，Kramers 定理保证每个时间反演不变动量上都有二重简并。

如果系统还具有空间反演 $P$，则
\[
P:\vb k\to-\vb k,
\qquad
\Theta:\vb k\to-\vb k,
\]
所以组合 $P\Theta$ 保持同一个 $\vb k$ 不变。对通常的非磁、中心对称自旋 $1/2$ 系统，$P$ 与 $\Theta$ 对易，且
\[
(P\Theta)^2=-1.
\]
于是 Kramers 型论证在\emph{每一个} $\vb k$ 上都成立，得到整条能带的双重简并。若破坏反演但保留时间反演，则一般 $\vb k$ 点可以发生自旋--轨道劈裂，只在时间反演不变动量处仍有 Kramers 简并。这就是 Rashba/Dresselhaus 型自旋分裂为何需要缺失空间反演的对称性基础。

时间反演对称下的 Kramers 二重简并能带如图 \figref{fig:tr-kramers-bands} 所示。
\begin{figure}[htbp]
\centering
\begin{tikzfig}[scale=0.98]
  \draw[axes] (-3.2,0)--(3.25,0) node[right] {$k$};
  \draw[axes] (0,-.4)--(0,4.6) node[above] {$E$};
  \draw[line width=.9pt,domain=-3:3,samples=80] plot (\x,{1.15+0.28*(\x)^2+0.42*\x});
  \draw[line width=.9pt,domain=-3:3,samples=80] plot (\x,{1.15+0.28*(\x)^2-0.42*\x});
  \fill (0,1.15) circle (1.6pt);
  \node[lecturelabel] at (0.85,0.55) {Kramers point};
  \node[lecturelabel] at (1.65,3.7) {$E_+(k)=E_-(-k)$};
\end{tikzfig}
\caption{保时间反演但缺少空间反演时的典型自旋分裂示意。一般 $k$ 处两支分开，但在时间反演不变动量处仍必须相交成 Kramers 对。}
\label{fig:tr-kramers-bands}
\end{figure}

时间反演的反酉性还解释了一个在点群特征标表中常见的现象。若一个无自旋态按复表示 $D(g)$ 变换，则其时间反演态中的系数被复共轭，因此按 $D(g)^*$ 变换。若 $D$ 与 $D^*$ 不等价，纯空间点群并不要求两者属于同一个不可约表示；但时间反演会把它们联系起来。在具有时间反演对称性的完整问题中，这对复共轭表示不能被任意分开处理。

这也说明为什么仅看普通点群特征标表有时会低估简并：真正的对称结构可能包含反酉操作。要系统处理磁性系统和时间反演，需要把普通群表示推广到共表示（corepresentation）理论；这超出本课程主体范围，但本节的 Kramers 定理和 $\vb k\leftrightarrow-\vb k$ 关系已经覆盖后续凝聚态物理中最常用的部分。


时间反演算符把复表示 $D(g)$ 送到复共轭 $D(g)^*$。一个自然的问题是：$D$ 与 $D^*$ 是否等价？如果等价，能否在实数域上实现？这两个问题的答案由\textbf{Frobenius--Schur 指标}完全分类。这个指标不仅是一个纯代数不变量，还直接决定时间反演对称性下简并的结构。

\begin{definitionplain}[Frobenius--Schur 指标]
设 $G$ 是有限群，$D^{(\alpha)}$ 是不可约酉表示，维数 $d_\alpha$。定义
\begin{equation}
\nu_\alpha
=\frac{1}{|G|}\sum_{g\in G}\chi^{(\alpha)}(g^2).
\label{eq:frobenius-schur-indicator}
\end{equation}
称为不可约表示 $\alpha$ 的\textbf{Frobenius--Schur 指标}。
\label{def:frobenius-schur-indicator}
\end{definitionplain}

\begin{theorem}[Frobenius--Schur 分类定理]
对有限群 $G$ 的每个不可约表示 $\alpha$，Frobenius--Schur 指标 $\nu_\alpha\in\{+1,0,-1\}$，且：
\begin{enumerate}[label=(\roman*),leftmargin=3em]
\item $\nu_\alpha=+1$（\textbf{实型}）：$D^{(\alpha)}\cong D^{(\alpha)*}$，且存在对称非退化不变双线性形式。表示可在实数域上实现。
\item $\nu_\alpha=0$（\textbf{复型}）：$D^{(\alpha)}\not\cong D^{(\alpha)*}$。$D^{(\alpha)}$ 与 $D^{(\alpha)*}$ 是两个不等价不可约表示。
\item $\nu_\alpha=-1$（\textbf{四元数型/伪实型}）：$D^{(\alpha)}\cong D^{(\alpha)*}$，且存在反对称非退化不变双线性形式。表示不能在实数域上实现。
\end{enumerate}
\label{thm:frobenius-schur-classification}
\end{theorem}

\par\medskip
\noindent\textbf{Frobenius--Schur 分类定理的证明}\quad
在 $V_\alpha\otimes V_\alpha$ 上定义群平均投影与交换算符
\[
P=\frac1{|G|}\sum_{g\in G}D^{(\alpha)}(g)\otimes D^{(\alpha)}(g),
\qquad \tau(v\otimes w)=w\otimes v.
\]
$P$ 投影到不变张量子空间，而恒等式
$\tr[\tau(A\otimes B)]=\tr(AB)$ 给出
\[
\tr(\tau P)
=\frac1{|G|}\sum_{g\in G}\tr D^{(\alpha)}(g)^2
=\frac1{|G|}\sum_{g\in G}\chi^{(\alpha)}(g^2)
=\nu_\alpha.
\]
不变张量等价于从 $V_\alpha^*$ 到 $V_\alpha$ 的交织映射。由 Schur 引理，
$V_\alpha\not\cong V_\alpha^*$ 时不变子空间为零，故 $\nu_\alpha=0$；自对偶时
该空间是一维。交换算符 $\tau$ 在这个一维空间上的本征值只能是 $+1$ 或 $-1$，
分别表示不变双线性形式为对称或反对称，因此 $\nu_\alpha=+1$ 或 $-1$。

在酉表示中，对称的非退化不变形式允许选择实正交基，所以 $+1$ 对应实型。
反对称非退化形式只存在于偶数维空间，并给出四元数结构，所以 $-1$ 对应
四元数型。三种情形由此穷尽。
\par\medskip

Frobenius--Schur 指标与时间反演的关系是本节的核心物理连接。

\begin{theorem}[Frobenius--Schur 指标与时间反演简并]
设反酉时间反演 $\Theta$ 与群 $G$ 的作用对易，并满足 $\Theta^2=-1$。若能级的空间对称类型为不可约表示 $\alpha$，则其最小时间反演闭合空间的维数为
\[
d_{\min}=\begin{cases}
d_\alpha,&\nu_\alpha=-1,\\
2d_\alpha,&\nu_\alpha=0\text{ 或 }+1.
\end{cases}
\]
四元数型表示内部已经包含 Kramers 配对；复型表示须与共轭表示 $\alpha^*$ 配对，实型表示则须复制一份才能实现 $\Theta^2=-1$。
\label{thm:fs-time-reversal-degeneracy}
\end{theorem}

\par\medskip
\noindent\textbf{定理证明}\quad
在不可约子空间 $V_\alpha$ 上，$\Theta$ 给出从 $D^{(\alpha)}$ 到 $D^{(\alpha)*}$ 的反线性交织。若 $\nu_\alpha=0$，两表示不等价，故 $\Theta V_\alpha$ 是另一个同能的 $V_{\alpha^*}$，最小闭合空间为 $V_\alpha\oplus V_{\alpha^*}$。

若 $\nu_\alpha=+1$，可取基使 $D^{(\alpha)}(g)$ 为实矩阵。在单个副本上，唯一的反线性交织在相位意义下是复共轭 $K$，其平方为 $+1$；要得到 $\Theta^2=-1$，必须在两个副本上取 $\Theta=(\ii\sigma_y)K$，故维数加倍。

若 $\nu_\alpha=-1$，存在非退化反对称不变矩阵 $J$，可归一化为 $JJ^*=-I$。于是 $\Theta=JK$ 已在 $V_\alpha$ 内满足 $\Theta^2=-1$。反对称非退化性还保证 $d_\alpha$ 为偶数，所以 Kramers 对已经包含在该不可约表示内部，无须再复制表示空间。
\par\medskip

在晶场问题中，这一区分说明了为何仅知道空间群不可约表示的维数还不够：半整数自旋还必须结合双群表示及其 Frobenius--Schur 类型，才能读出实际简并度。

\medskip
\noindent\emph{$SU(2)$ 表示的 Frobenius--Schur 指标。}
对 $SU(2)$ 的 $j$ 表示（维数 $2j+1$），特征标为 $\chi_j(\theta)=\frac{\sin((2j+1)\theta/2)}{\sin(\theta/2)}$。Frobenius--Schur 指标为
\[
\nu_j=(-1)^{2j}.
\]
因此：
\begin{itemize}
\item $j\in\NN$（整数自旋）：$\nu=+1$，实型。表示可在实数域上实现（如 $j=1$ 的三维表示是实正交矩阵）。
\item $j\in\NN+1/2$（半整数自旋）：$\nu=-1$，四元数型。表示维数 $2j+1$ 为偶数，存在反对称不变形式。这正是 Kramers 简并的群论根源。
\end{itemize}
对 $SU(2)$ 没有复型（$\nu=0$）表示，因为所有表示都自共轭：$D^{(j)*}\cong D^{(j)}$。这反映了 $SU(2)$ 的特殊性——其所有不可约表示都是实或四元数型。
\par\medskip

\medskip
\noindent\emph{点群的 Frobenius--Schur 指标。}
对常见晶体点群：
\begin{itemize}
\item $C_n$ 的一维表示 $e^{2\pi im/n}$：若 $m=0$ 或 $m=n/2$（当 $n$ 偶），则 $\nu=+1$（实型）；否则 $\nu=0$（复型），因为 $e^{2\pi im/n}$ 与 $e^{-2\pi im/n}$ 不等价。
\item $D_n$ 的所有不可约表示都是实型（$\nu=+1$），因为 $D_n$ 是实反射群。
\item $T$（正四面体群）的三维表示：$\nu=+1$（实型，由实正交矩阵实现）。
\item $O$（立方群）的三维表示：$\nu=+1$。
\end{itemize}
晶体点群中四元数型表示（$\nu=-1$）很少见，但双群中会出现。例如 $T_d$ 的双群 $\bar T_d$ 有二维旋量表示，其 $\nu=-1$，这正是半整数自旋在晶场中的 Kramers 简并来源。
\par\medskip

Frobenius--Schur 指标还给出一个有用的维数约束。

\begin{proposition}[四元数型表示维数为偶数]
若 $\nu_\alpha=-1$，则 $d_\alpha$ 为偶数。
\label{prop:quaternionic-even-dim}

\end{proposition}

\begin{proof}
四元数型表示存在非退化反对称不变双线性形式 $B$。非退化反对称矩阵只存在于偶数维空间（$\det A=\det A^T=(-1)^{d_\alpha}\det A$，若 $d_\alpha$ 奇则 $\det A=0$）。故 $d_\alpha$ 为偶数。
\end{proof}

这个结论与 Kramers 定理一致：半整数自旋的不可约表示维数 $2j+1$ 为偶数，因此可以承载 $\Theta^2=-1$ 的反酉时间反演。

\subsection{对称性破缺与序参量}
\label{sec:symmetry-breaking}

前面各节处理的是"对称性 $\Rightarrow$ 简并与选择定则"的正向逻辑。现在反过来：当系统的对称性从高对称群 $G$ 降为子群 $H\subset G$ 时，群论如何描述这一"破缺"？这个问题在相变理论（Landau 理论）、Jahn--Teller 效应和凝聚态物理中是核心。


\begin{definitionplain}[对称性破缺]
设系统 Hamilton 量 $\hat H$ 在群 $G$ 下不变（$[\hat H,U(g)]=0$，$\forall g\in G$）。若基态 $\ket{\psi_0}$ 满足 $U(g)\ket{\psi_0}\ne\ket{\psi_0}$ 对某些 $g\in G$，则称对称性从 $G$ 破缺到\textbf{保持基态不变的子群}
\[
H=\{h\in G\mid U(h)\ket{\psi_0}=\ket{\psi_0}\}\le G.
\]
$H$ 称为\textbf{残余对称群}，$G/H$（若 $H\mathrel{\trianglelefteq} G$）或陪集结构描述被破缺的对称操作。
\end{definitionplain}

\begin{proposition}[序参量的表示论]
\label{prop:order-parameter}
对称性破缺由\textbf{序参量} $\eta$ 描述，它是在 $G$ 下按某个非平凡不可约表示 $\Gamma\ne A_1$ 变换的量。具体地：
\begin{enumerate}[label=(\roman*),leftmargin=3em]
\item 高对称相：$\eta=0$，系统在 $G$ 下不变。
\item 低对称相：$\eta\ne0$，$\eta$ 在残余群 $H$ 下不变，但在 $G\setminus H$ 的操作下改变。
\item $\eta$ 按 $\Gamma$ 的某个分量变换，$H$ 是 $\Gamma$ 中使该分量不变的子群。
\end{enumerate}
\end{proposition}

\begin{proof}
(i) 高对称相中基态 $\ket{\psi_0}$ 在 $G$ 下不变，所有按非平凡表示变换的量期望值为零（由 \cref{thm:matrix-element-selection}，$A_1\to\Gamma$ 的矩阵元为零若 $\Gamma\ne A_1$）。

(ii) 低对称相中 $\ket{\psi_0}$ 只在 $H$ 下不变。序参量 $\eta=\bra{\psi_0}\hat O\ket{\psi_0}$（$\hat O$ 按 $\Gamma$ 变换）非零，因 $\hat O$ 在 $H$ 下有 $A_1$ 分量（即 $\Gamma\downarrow_H$ 含 $A_1$）。但对 $g\in G\setminus H$，$U(g)\ket{\psi_0}\ne\ket{\psi_0}$，$\eta$ 在 $g$ 下改变。

(iii) $H=\{g\in G\mid D^{(\Gamma)}(g)\eta=\eta\}$，即 $H$ 是 $\eta$ 在 $\Gamma$ 表示中的稳定子群。
\end{proof}

\medskip
\noindent\emph{铁磁相变的序参量。}
各向同性铁磁体的对称群 $G=SO(3)$（自旋转动）。高温顺磁相：$\vb M=0$（磁化强度为零），$SO(3)$ 完整。低温铁磁相：$\vb M\ne0$（自发磁化），设 $\vb M$ 沿 $z$ 方向。残余对称群 $H=SO(2)$（绕 $z$ 轴的转动保持 $\vb M$ 不变）。序参量 $\vb M$ 是矢量（秩 $1$ 张量），按 $SO(3)$ 的 $j=1$ 表示变换。破缺的对称操作：$SO(3)/SO(2)\cong S^2$（单位球面），对应磁化方向的选择——这就是铁磁体中"自发选择方向"的群论描述。
\par\medskip


Landau 相变理论把自由能展开为序参量的幂级数，群论严格约束哪些项可以出现。

\begin{theorem}[Landau 自由能的群论约束]
\label{thm:landau-group}
设序参量 $\eta$ 按群 $G$ 的不可约表示 $\Gamma$ 变换（维数 $d_\Gamma$）。自由能 $F(\eta)$ 在 $G$ 下不变，故只含 $G$ 不变的项。具体地：
\begin{enumerate}[label=(\roman*),leftmargin=3em]
\item 线性项 $\eta_i$：仅当 $\Gamma=A_1$（平凡表示）时允许。对非平凡 $\Gamma$，$\sum_i\eta_i$ 不是不变量，线性项被禁。
\item 二次项 $\eta_i\eta_j$：对称化后按 $\mathrm{Sym}^2\Gamma$ 变换，不变项要求 $A_1\subset\mathrm{Sym}^2\Gamma$。对实表示，$A_1$ 总在 $\mathrm{Sym}^2\Gamma$ 中（$\sum_i\eta_i^2$ 不变），给出 $\alpha|\eta|^2$ 项。
\item 三次项 $\eta_i\eta_j\eta_k$：按 $\mathrm{Sym}^3\Gamma$ 变换，仅当 $A_1\subset\mathrm{Sym}^3\Gamma$ 时允许。
\item 四次项：按 $\mathrm{Sym}^4\Gamma$ 变换，$A_1\subset\mathrm{Sym}^4\Gamma$ 的重数给出独立四次不变量个数。
\end{enumerate}
\end{theorem}

\begin{proof}
自由能 $F$ 是标量（在 $G$ 下不变），故 $F$ 中每项必须是 $G$ 不变的多项式。秩 $n$ 不变量对应 $\mathrm{Sym}^n\Gamma$ 中的 $A_1$ 分量，重数 $m_{A_1}=\frac1{|G|}\sum_g\chi_{\mathrm{Sym}^n\Gamma}(g)$ 给出独立的 $n$ 次不变量个数。

线性项。$\eta_i$ 按 $\Gamma$ 变换，$A_1\subset\Gamma$ 当且仅当 $\Gamma=A_1$。对非平凡序参量，线性项被禁——这就是 Landau 理论中"连续相变要求序参量在临界点连续变化"的群论基础。

二次项。$\eta_i\eta_j$ 按 $\Gamma\otimes\Gamma$ 变换，对称部分按 $\mathrm{Sym}^2\Gamma$ 变换。对实表示，$\Gamma\otimes\Gamma$ 总含 $A_1$（因 $\chi_\Gamma^2$ 的平均 $\frac1{|G|}\sum|\chi_\Gamma|^2=1$），且 $A_1$ 在对称部分（因 $\sum_i\eta_i^2$ 不变）。故二次项 $\alpha\sum_i\eta_i^2$ 总允许。

三次项。$\mathrm{Sym}^3\Gamma$ 是否含 $A_1$ 取决于 $\Gamma$ 的具体结构。若 $A_1\not\subset\mathrm{Sym}^3\Gamma$，三次项被禁，Landau 自由能 $F=\alpha|\eta|^2+\beta|\eta|^4+\cdots$ 允许连续的二阶相变。若 $A_1\subset\mathrm{Sym}^3\Gamma$，三次项允许，一般导致一阶相变（序参量在临界点跳变）。
\end{proof}

\medskip
\noindent\emph{Landau 理论对立方晶系铁电相变的应用。}
立方晶系铁电体（如 $\mathrm{BaTiO_3}$）的序参量是极化矢量 $\vb P=(P_x,P_y,P_z)$，按 $O$ 的三维不可约表示 $T_1$ 变换。

二次项。$\mathrm{Sym}^2 T_1=A_1\oplus E\oplus T_2$，$A_1$ 重数 $1$，唯一二次不变量 $P_x^2+P_y^2+P_z^2=|\vb P|^2$。

三次项。$\mathrm{Sym}^3 T_1=A_1\oplus A_2\oplus 2E\oplus 2T_1\oplus T_2$，$A_1$ 重数 $1$，唯一三次不变量 $P_xP_yP_z$。三次项允许意味着立方铁电相变一般是一阶的——这正是 $\mathrm{BaTiO_3}$ 相变为一阶的群论原因。

四次项。$\mathrm{Sym}^4 T_1$ 中 $A_1$ 重数 $2$，两个独立四次不变量：$|\vb P|^4=(P_x^2+P_y^2+P_z^2)^2$ 和 $P_x^4+P_y^4+P_z^4$。自由能
\[
F=\alpha|\vb P|^2+\beta P_xP_yP_z+\gamma_1|\vb P|^4+\gamma_2(P_x^4+P_y^4+P_z^4)+\cdots.
\]
$\gamma_2$ 项决定极化方向：$\gamma_2<0$ 时沿 $\langle 100\rangle$（如四方 $\mathrm{BaTiO_3}$ 低温相），$\gamma_2>0$ 时沿 $\langle 111\rangle$（如三方 $\mathrm{LiNbO_3}$）。
\par\medskip


Jahn--Teller 效应是分子和晶体中对称性自发破缺的经典实例：高对称构型下若电子态简并，系统会自发畸变以消除简并，降低能量。

\begin{theorem}[Jahn--Teller 定理]
\label{thm:jahn-teller}
设分子的电子基态在核构型 $Q_0$（高对称点群 $G$）下按不可约表示 $\Gamma$ 变换，$\dim\Gamma>1$（简并）。若振动表示 $\rho_{\mathrm{vib}}$ 的对称化平方 $\mathrm{Sym}^2\Gamma$ 与某个非全对称振动模式 $\Gamma_v\ne A_1$ 有非零交集：
\[
\mathrm{Sym}^2\Gamma\otimes\Gamma_v\supset A_1,
\]
则 $Q_0$ 不是稳定平衡构型：存在畸变 $Q\ne Q_0$ 使能量降低，对称性从 $G$ 破缺到 $H\subset G$。
\end{theorem}

\begin{proof}
先能量展开。把电子能量 $E(Q)$ 在 $Q_0$ 附近按核位移 $Q$ 展开：
\[
E(Q)=E_0+\sum_i\frac{\partial E}{\partial Q_i}Q_i+\frac12\sum_{ij}\frac{\partial^2E}{\partial Q_i\partial Q_j}Q_iQ_j+\cdots.
\]
高对称点 $Q_0$ 是平衡点，一次项为零。

随后线性耦合。简并态 $\{\ket{\Gamma,m}\}$ 的能量矩阵为
\[
H_{mn}(Q)=E_0\delta_{mn}+\sum_i\bra{\Gamma,m}\frac{\partial\hat H}{\partial Q_i}\ket{\Gamma,n}Q_i+\cdots.
\]
矩阵元 $\bra{\Gamma,m}\partial\hat H/\partial Q_i\ket{\Gamma,n}$ 非零要求 $A_1\subset\Gamma^*\otimes\Gamma_v\otimes\Gamma$（$Q_i$ 按 $\Gamma_v$ 变换，$\partial\hat H/\partial Q_i$ 也按 $\Gamma_v$ 变换）。等价地 $\Gamma\otimes\Gamma\otimes\Gamma_v\supset A_1$，即 $\mathrm{Sym}^2\Gamma\otimes\Gamma_v\supset A_1$（因 $m,n$ 对称）。

接着能量降低。若条件满足，线性耦合矩阵 $V_{mn}=\sum_i\bra{\Gamma,m}\partial\hat H/\partial Q_i\ket{\Gamma,n}Q_i$ 非零。$V$ 是 $\dim\Gamma\times\dim\Gamma$ 厄米矩阵，迹为零（因 $V$ 按 $\Gamma_v\ne A_1$ 变换，$A_1$ 分量的迹为零）。迹零厄米矩阵至少有一个负本征值。故存在简并态的线性组合使能量降低 $\Delta E<0$，畸变 $Q$ 使总能量 $E(Q)<E_0$。

再对称性破缺。畸变 $Q$ 按 $\Gamma_v\ne A_1$ 变换，故 $Q$ 不在所有 $G$ 操作下不变。残余对称群 $H=\{g\in G\mid D^{(\Gamma_v)}(g)Q=Q\}$ 是 $G$ 的真子群。
\end{proof}

\begin{exampleplain}[$E\otimes e$ Jahn--Teller 效应]
正三角形分子（如 $\mathrm{H_3}$、$\mathrm{C_3H_3}$）的电子基态按 $C_{3v}$ 的 $E$ 表示变换（二重简并），振动模式中有一个 $e$ 型模式（二重简并的弯曲振动）。检查条件：
\[
\mathrm{Sym}^2 E = A_1\oplus E,\qquad \mathrm{Sym}^2 E\otimes E = (A_1\oplus E)\otimes E = E\oplus(A_1\oplus A_2\oplus E)\supset A_1.
\]
条件满足，Jahn--Teller 效应发生：正三角形构型不稳定，系统畸变为等腰三角形（对称性从 $C_{3v}$ 破缺到 $C_s$）。这就是 $E\otimes e$ Jahn--Teller 效应，是分子物理中最常见的类型。

例外。若 $\Gamma$ 是一维表示（$\dim\Gamma=1$），$\mathrm{Sym}^2\Gamma=\Gamma^2$，对一维表示 $\Gamma^2$ 可能不含任何 $\Gamma_v\ne A_1$，Jahn--Teller 效应不发生。线性分子的 $\Sigma$ 态（一维）不发生 Jahn--Teller 畸变，保持线性构型。Kramers 简并（时间反演强制二重简并）也不受 Jahn--Teller 效应影响——Kramers 对不能被晶场劈裂，这是 \cref{thm:kramers} 的推论。
\end{exampleplain}

\subsection{习题与解答}
\label{sec:quantum-group-exercises}

下面的习题依次处理本征空间上的表示、对称性劈裂、投影算符、选择定则、Bloch 定理与时间反演。每题解答都明确指出所调用的定义或定理。

\begin{enumerate}[label=\arabic*.,leftmargin=2.5em]

\item 设 $[\hat H,U(g)]=0$。证明若 $\ket\psi$ 是能量 $E$ 的本征态，则 $U(g)\ket\psi$ 也是能量 $E$ 的本征态。

\emph{解答.}
\[
\hat H U(g)\ket\psi
=U(g)\hat H\ket\psi
=E U(g)\ket\psi.
\]
这正是 \Cref{thm:eigenspace-invariant} 的单态版本。

\item 某能级承载点群的一个三维不可约表示。只由该点群对称性，能否保证该能级恰好三重简并？

\emph{解答.} 只能保证\emph{至少}三重简并。若同一能量还包含其他不可约表示，或者同一三维表示出现多份并且由于额外动力学原因碰巧同能，简并度可以更高。对称性强制的是不可约表示维数这一基本单位。

\item 设原对称群 $G$ 的二维不可约表示在子群 $H$ 下限制为 $A\oplus B$，其中 $A,B$ 是两个不等价一维表示。说明一个保持 $H$ 但破坏 $G$ 的一般微扰会怎样影响原二重简并。

\emph{解答.} 根据 \Cref{thm:perturbation-block-structure}，一阶微扰矩阵在 $A$ 与 $B$ 之间没有矩阵元，各自形成一个一维块。由于 $A$、$B$ 不等价，残余对称性不要求两个块能量相同，因此一般会分裂成两个非简并能级。

\item 写出 $O_h$ 晶场中五个 $d$ 轨道的不可约分解，并给出一组对应基函数。

\emph{解答.}
\[
\Gamma_d=E_g\oplus T_{2g},
\]
其中
\[
E_g:\ \{2z^2-x^2-y^2,\ x^2-y^2\},
\qquad
T_{2g}:\ \{xy,yz,zx\}.
\]

\item 若把上一题的立方晶场沿 $z$ 轴作四方畸变，使对称性降为 $D_{4h}$，写出允许的进一步分裂。

\emph{解答.}
\[
E_g(O_h)\to A_{1g}\oplus B_{1g},
\qquad
T_{2g}(O_h)\to B_{2g}\oplus E_g(D_{4h}).
\]
因此五重空间最终按 $1+1+1+2$ 分组。

\item 对\figref{fig:c3v-three-sites} 的三个等价位置，证明全对称态
\[
\ket{A_1}=\frac1{\sqrt3}(\ket1+\ket2+\ket3)
\]
在所有 $C_{3v}$ 操作下不变。

\emph{解答.} $C_{3v}$ 的任何群元只会置换 $\ket1,\ket2,\ket3$。三者系数完全相同，因此置换后线性组合不变，即承载平凡不可约表示 $A_1$。

\item 设 $C_{3v}$ 对称哈密顿量在三个等价位置基下为
\[
H=\begin{pmatrix}\varepsilon&t&t\\t&\varepsilon&t\\t&t&\varepsilon\end{pmatrix}.
\]
求全部本征值，并指出对应表示。

\emph{解答.} 全对称态给
\[
E_{A_1}=\varepsilon+2t.
\]
与它正交的二维子空间是 $E$ 表示，任一向量都满足分量和为零，因此
\[
H\vb v=(\varepsilon-t)\vb v.
\]
故
\[
E_E=\varepsilon-t
\]
二重简并。

\item 在中心对称体系中，为什么电偶极跃迁要求初末态宇称相反？

\emph{解答.} 电偶极算符 $\hat{\vb r}$ 在反演下为奇，属于 $u$。非零矩阵元要求
\[
\Gamma_f^*\otimes u\otimes\Gamma_i
\]
含全对称 $g$。因此 $\Gamma_f$ 与 $\Gamma_i$ 的宇称必须相反，即 $g\leftrightarrow u$。

\item 在 $C_{4v}$ 中，初态属于 $A_1$。分别判断 $z$ 偏振与 $x$ 偏振电偶极光能把它耦合到哪些不可约表示。

\emph{解答.} $z\sim A_1$，故
\[
A_1\otimes A_1=A_1,
\]
所以 $z$ 偏振允许到 $A_1$ 末态。$(x,y)\sim E$，故 $x$ 或 $y$ 偏振可把 $A_1$ 耦合到 $E$ 末态。

\item 某中心对称分子的振动模为 $g$ 型。它能否同时是一阶电偶极 IR 主动与普通 Raman 主动？

\emph{解答.} 不能。IR 偶极算符为 $u$，因此 IR 主动模必须为 $u$；Raman 张量由两个矢量乘积组成，为 $g$，所以 $g$ 型模可以 Raman 主动。中心对称体系中两种活性互斥。

\item 证明体中心对称晶体的电偶极二阶极化率 $\chi^{(2)}$ 必为零。

\emph{解答.} 反演下 $\vb P\to-\vb P$，而 $E_jE_k$ 为偶。若
\[
P_i^{(2)}=\epsilon_0\chi^{(2)}_{ijk}E_jE_k
\]
在反演后形式不变，则必须同时满足
\[
-P_i^{(2)}=\epsilon_0\chi^{(2)}_{ijk}E_jE_k=P_i^{(2)},
\]
故所有电偶极 $\chi^{(2)}_{ijk}=0$。

\item 由 $T_{\vb R}\psi_{\vb k}=\ee^{-\ii\vb k\cdot\vb R}\psi_{\vb k}$ 推出 Bloch 形式。

\emph{解答.} 由平移算符定义得到
\[
\psi_{\vb k}(\vb r+\vb R)=\ee^{\ii\vb k\cdot\vb R}\psi_{\vb k}(\vb r).
\]
定义
\[
u_{\vb k}(\vb r)=\ee^{-\ii\vb k\cdot\vb r}\psi_{\vb k}(\vb r),
\]
则直接代入可得 $u_{\vb k}(\vb r+\vb R)=u_{\vb k}(\vb r)$，故
\[
\psi_{\vb k}=\ee^{\ii\vb k\cdot\vb r}u_{\vb k}.
\]

\item 为什么 $\vb k$ 与 $\vb k+\vb G$ 标记相同的平移群不可约表示？

\emph{解答.} 对任意晶格矢量 $\vb R$，
\[
\ee^{-\ii(\vb k+\vb G)\cdot\vb R}
=\ee^{-\ii\vb k\cdot\vb R}\ee^{-\ii\vb G\cdot\vb R}
=\ee^{-\ii\vb k\cdot\vb R},
\]
因为 $\vb G\cdot\vb R\in2\pi\mathbb Z$。

\item 若空间群中存在二重螺旋轴 $S=\{C_{2z}|c/2\}$，证明在 $k_z=\pi/c$ 上其本征值为 $\pm\ii$。

\emph{解答.} $S^2=T_c$，因此
\[
s^2=\ee^{-\ii k_z c}.
\]
当 $k_z=\pi/c$ 时，$s^2=-1$，所以
\[
s=\pm\ii.
\]

\item 对自旋 $1/2$，验证 $\Theta=-\ii\sigma_yK$ 满足 $\Theta^2=-1$。

\emph{解答.} 注意 $K(-\ii\sigma_y)K^{-1}=(-\ii\sigma_y)^*$。由于 $-\ii\sigma_y$ 是实矩阵
\[
-\ii\sigma_y=\begin{pmatrix}0&-1\\1&0\end{pmatrix},
\]
故
\[
\Theta^2=(-\ii\sigma_y)^2=-I.
\]

\item 证明具有时间反演和空间反演的自旋 $1/2$ 非磁晶体在每个 $\vb k$ 上至少二重简并。

\emph{解答.} $P$ 与 $\Theta$ 都把 $\vb k$ 送到 $-\vb k$，所以 $P\Theta$ 保持 $\vb k$。若二者对易，则
\[
(P\Theta)^2=P^2\Theta^2=-1.
\]
同时 $P\Theta$ 与哈密顿量对易，因此对每个固定 $\vb k$ 都可应用 Kramers 论证，得到二重简并。

\end{enumerate}
